\begin{table}[htbp]
  \centering\small
  \caption{31 场景缺失率按类型汇总（附件2 test）}
  \label{tab:p2-missrate-summary}
  \begin{tabular}{l r c c c c}
    \toprule
    场景类型 & 场景数 & 均值 & 标准差 & 最小 & 最大 \\
    \midrule
    MCAR 单模态（$\{t,a,v\}\times\{10,20,40,60,80\}\%$）   & 15 & 0.4201 & 0.2643 & 0.1004 & 0.7985 \\
    MCAR 双模态（$\{ta,tv,av\}@40\%$）                     &  3 & 0.3999 & 0.0000 & 0.3999 & 0.3999 \\
    连续块（$\{t,a,v\}\times\{$头/中/尾/随机$\}@40\%$）   & 12 & 0.3976 & 0.0037 & 0.3889 & 0.3999 \\
    联合缺失（$av@40\%$）                                  &  1 & 0.3846 & ---    & 0.3846 & 0.3846 \\
    \midrule
    31 场景合计                                            & 31 & 0.4074 & 0.2179 & 0.1004 & 0.7985 \\
    \bottomrule
  \end{tabular}
  \par \smallskip
  {\small 注：与图~\ref{fig:dist-miss-rate}~共同支撑——三划分密度曲线高度重叠，证明网格覆盖率分布在训练/验证/测试同源；连续块场景实际缺失率略低于名义 40\%（最差 v 系列 38.9\%），源于连续块在头/中/尾受限的可铺区域。联合缺失组 n=1 无标准差。}
\end{table}